% !TEX root = ../main.tex
% Kapitel 1: Einleitung und Motivation
% Quelle/Ausgangspunkt: Projektanweisung_Invertiertes_Pendel.md (Abschnitte 1-2), eigene Motivation.
\chapter{Einleitung und Motivation}
\label{ch:einleitung}
\section{Problemstellung und Benchmark-Charakter}
Das mechatornische System des invertierten Pendels in diesem Projekt besteht aus einem translatorisch gelagerten
Wagen auf einer horizontalen Schiene und einem daran befestigten Pendelstab, auf diesem Stab wirkt keine direkte Stellgröße.
Dieses System bestizt in der Regelungstechnik und der nichtlinearen Dynamik eine besondere Benchmark-Charakter zur Validierung von Regelungsalgorithmen/ mathematische Kontrollstrategien (vgl. Stand der Technik des invertierten Pendels und dessen Integration in das Python-Software-Ökosystem
). Die Komplexität der Stabilisierung des invertierten Pendels auf dem Wagen resultiert aus drei physikalischen Eigenschaften:
\begin{itemize}
    \item Das System ist unterbestimmt, da die zwei Freiheitsgrade (Wagenposition s und Pendelwinkel $\phi$) nur durch eine Stellgröße $\tau$ die auf den Wagen wirkt beeinflusst werden können.
    \item Die Dynamik des Systems ist nichtlinear, da die Bewegungsgleichungen des invertierten Pendels nichtlinear sind.
    \item Die Zielposition (Pendel aufrecht, $\phi = \tau$) ist ein instabiler Gleichgewichtszustand, bei der kleinste Störungen zu einem Umkippen des Pendels führen können.
\end{itemize}
Einen ersten visuellen Eindruck der Simulationsumgebung ist in Abbildung \ref{fig:pendulum_simulation_interactive} dargestellt.
\begin{figure}[htbp]
    \centering
    \includegraphics[width=0.75\textwidth]{game_interactive}
    \caption{Interaktive Simulationsumgebung des invertierten Pendels (Pygame-Fenster im laufenden Spiel).}
    \label{fig:pendulum_simulation_interactive}
\end{figure}

\section{Stand der Technik}
% mache ich später fertig

\section{Zielsetzung und Ausgangslage}
Das Ziel des Projekts ist die Weiterentwicklung eines unvollständigen, rein formalbasierten OpenModelica-Referenzmodells (pendulum.mo) zu einem physikalisch realistischen, interaktiven und spielbaren Simulationsumgebung. Darauf hin wurde auch die bestehnde Python-Simulationumgebung (pendulum\_game\_controlled.py) erweitert. 
Die Ausgangslage vor Bedinn des Projekts umfasste ein flaches Modelica-Modell des invertierten Pendels, das die Bewegungsgleichungen des Systems formal korrekt beschreibt, jedoch keine physikalisch realistische Simulation ermöglicht.
Im Rahmen dieses Projekts wurde das bestehende Referenzmodell durch eine objektorientierte Implementierung unter Verwendung der Modelica Standard Library (MSL), insbesondere der MultiBody-Bibliothek ersetzt um die gekoppelte Dynamik abzubilden.
Um die Interaktivität und Spielbarkeit zu gewährleisten, wurde die kompilierte Simulations-FMU in eine Python-Umgebung eingebunden.
Diese wurde um Tastensteuerung (Pause- und Reset-Funktion), ein zonebasiertes Punktesystem mit kontinuierlichkeits Bonus sowie Schwierigkeitsgrade erweitert.

\section{Aufbau der Arbeit}
Die Ausarbeitung dokumentiert den Entwicklungsprozess und ist in die folgenden Kapitel gegliedert:
\begin{itemize}
    \item Kapitel 2 widmet sich den theoretischen Grundlagen der Modellierung und der Herleitung der Bewegungsgleichungen über die Euler-Lagrange-Gleichungen her.
    \item Kapitel 3 beschreibt den Aufbau des MultiBody-Modells in OpenModelica sowie dessen algebraische Validierung.
    \item Kapitel 4 befasst sich mit der softwaretechnischen Umsetung der Simulationsumgebung der Spielmechaniken und des zonebasierten Punktesystems.
    \item Kapitel 5 stellt den Entwurf und die Implementierung der automatisierten PD-, LQR- und Swing-Up-Regler vor, zudem wird die Performance der Regler in verglichen.
    \item Kapitel 6 analysiert die physikalischen und spielerischen Auswirkungen der Schwierigkeitsgrade sowie die behebung des Resonanzproblems beim Scoring.
    \item Die Erkenntisse aus den vorherigen Kapiteln und Sovler-Eigenschaften werden in Kaptiel 7 diskuiert, bevor die Arbeit in Kapitel 8 mit einem Fazit und einem Ausblick schließt. 
\end{itemize}
    % TODO: Motivation, Zielsetzung des Projekts, Aufbau der Arbeit.
